%%%%%%%%%%%%%%%%%%%%%%%%%%%%%%%%%%%%%%%%

\newcommand{\subhotelinf}{Hotel Infinity}
\subsection{\subhotelinf}
\label{sub_hotel_infinity}

We end this textbook as we began it, with a puzzle.

\begin{act}[Hotel Infinity Puzzle]
\label{act_hotel} 
~
    \begin{actpart}
        \item Watch Hotel Infinity, Part One at  \url{https://tinyurl.com/infhotel1}\footnote{videos on Hotel Infinity by mathematician Tova Brown posted to YouTube\texttrademark.}.  Pause at the end of Part One.  Hotel Infinity is full one new guest arrives, the President of Space. How do you find a room for her?  (Watch video for rules.)
        
        \item Watch Hotel Infinity, Part Two at \url{https://tinyurl.com/infhotel2}. Pause at the end of Part Two.  Hotel Infinity is full when a bus load of infinitely many football players shows up. How do you find rooms for them?

        \item Watch Hotel Infinity, Part Three at \url{https://tinyurl.com/infhotel3}. Pause at the end of Part Three.  Hotel Infinity is empty when an infinite number of busses, each with infinitely many football players show up.  How do you find rooms for them?

        \item Watch Hotel Infinity, Part Four at \url{https://tinyurl.com/infhotel4}. Pause at the end of Part Four.  Hotel Infinity is empty when Visitors arrive (in a blob).  How do you find rooms for them?

        \item Watch Hotel Infinity, Part Five at \url{https://tinyurl.com/infhotel5}. Pause at the end of Part Five. Why is it impossible to fit all the Visitors in Hotel Infinity?
    \end{actpart}
\end{act}



%%%%%%%%%%%%%%%%%%%%%%%%%%%%%%%%%%%%%%%%

\newcommand{\subdefndiscrete}{The Definition of Discrete}
\subsection{\subdefndiscrete}
\label{sub_defn_discrete}

We can formalize our findings about Hotel Infinity by defining what it means for (infinite) sets to be the same size.

\begin{defn}[Countable and Uncountable Sets]
\label{defn_countable}
~
    \begin{apart}
        \item Sets $S$ and $T$ \vocab{are the same size set (have the same cardinality)} if there exists a pairing between $S$ and $T$, as defined in \Cref{defn_pairing}. Equivalently, if there exists a function $f:S \to T$ that is a bijection (onto and one-to-one).  For example, $\{0,1,2,3\}$ is the same size as $\{5,6,7,8\}$ because we have the pairing $0 \leftrightarrow 5, 1 \leftrightarrow 6, 2 \leftrightarrow 7$, $3 \leftrightarrow 8$ or, equivalently the function $f:\{0,1,2,3\} \to \{5,6,7,8\}$ defined by $f(n)=n+5$ is onto and one-to-one.
        \item A set $T$  is \vocab{countable} if it is the same size as the natural numbers $\mathbb{N}$.  That is, if there is a bijection $f:\mathbb{N} \to T$ that is onto and one-to-one. For example, $T = \{n \in \mathbb{Z} \mid n \ge 5\} = \{5,6,7,8, \ldots\}$ is countable because $f:\mathbb{N} \to T$ defined by $f(n)=n+5$ is onto and one-to-one.
        \item A set $T$ is \vocab{uncountable} if it is not countable.  For example, the interval $[0,1] = \{x \in \mathbb{R} \mid 0 \le x \le 1\}$ is uncountable, as shown in the video Hotel Infinity, Part Five.
    \end{apart}
\end{defn}

There are several equivalent characterizations of countable.

\begin{rem}[Alternative characterizations of countable]
\label{rem_alt_defn_countable}
    The following statements are equivalent.
        \begin{apart}
            \item The set $T$ is countable. (There is a bijection $f: \mathbb{N} \to T$.)
            \item The set $T$ fits into and fills Hotel Infinity.
            \item \label{part_seq_countable} There is a sequence that lists each element of $T$ exactly once. (Recall that a function with domain $\mathbb{N}$ is a sequence starting at 0.)
        \end{apart}
\end{rem}

Let's show that the integers are countable using \Cref{rem_alt_defn_countable} (\ref{part_seq_countable}).

\begin{exam}[The integers are countable]
\label{exam_integers_countable}
    Show that the integers are countable by finding a sequence that lists each integer exactly once.
        \begin{soln}
            Consider the sequence $$0,1,-1,2,-2,3,-3,4,-4, \ldots.$$
            The pattern defining this sequence should be clear but we can write a rule for fun.  For $n \ge 0$, define
                $$a_n = \left\{  
                \begin{array}{cl}
                    0 & \text{if } n=0 \vspace{.05in}\\
                    \frac{n+1}{2} & \text{if } n\text{ is odd} \vspace{.05in}\\
                    ^-\frac{n}{2} & \text{if } n\text{ is even} 
                \end{array}
                \right.$$
        \end{soln}
\end{exam}





We are now ready to define the topic of this course.

\begin{defn}[Discrete]
\label{defn_discrete}
   A set is \vocab{discrete} if it is finite or countable.
\end{defn}

I think I'll stop here.\footnote{``I think I'll stop here,'' This is how, on 23rd of June 1993, Andrew Wiles ended his series of lectures at the Isaac Newton Institute in Cambridge. The applause, so witnesses report, was thunderous: Wiles had just delivered a proof of a result that had haunted mathematicians for over 350 years: Fermat's last theorem. Source: \url{https://wild.maths.org/andrew-wiles-and-fermats-last-theorem}}






% %%%%%%%%%%%%%%%%%%%%%%%%

% \subsection*{Exercises}

% \begin{exer}[\exerP Hotel Infinity]
% \label{exer_hotel}
%     You are the proprietor of Hotel Infinity. Your hotel has rooms numbered $1, 2, 3, 4, 5, \ldots$ -- one room for each positive integer.  Sometimes guests arrive to your hotel in special buses.  Each \textbf{bus} holds passengers numbered $1, 2, 3, 4, 5, \ldots$ -- one passenger for each positive integer.  You may ask guests to move rooms, but never more than once a day. Guests are never asked to share a room with anyone else.  If at all possible, you won't turn any guest away.
%         \begin{exerpart}
%             \item On Monday your hotel was empty when four soccer players arrived so you put them in rooms 1, 2, 3, and 4.  But then a bus pulls up.  What should you do? Identify the room number for person $n$ from the bus.
%             \item On Tuesday all your guests are still there when  a famous chef show up needing a room. How can you make room for the chef?  Be specific.
%             \item Wednesday, the four soccer players leave.  You likes to boast a 100\% occupancy rate so you'd like to move the guests so all the rooms are full again.  How can you do that? Hint: where were the soccer players?    
%         \end{exerpart}  
% \end{exer}

% \begin{exer}[\exerU Hotel Infinity with buses]
% \label{exer_hotel_buses}
%     You are the proprietor of Hotel Infinity, as in \Cref{exer_hotel}. 
%         \begin{exerpart}
%             \item Your empty us again when two buses show up at the same time.  How can you assign rooms? 
%             \item The next day the hotel is empty again when three buses show up.  How can you assign rooms?
%             \item The next day the hotel is empty again and the buses just keep coming \ldots Bus \#1, Bus \#2, Bus \#3, Bus \#4, \ldots -- one bus for each positive integer.  Can you find a room for each guest?  If so, how?  Hint:  explain how to line up all the guests.  No $\ldots$ allowed in your line-up, except at the very end. 
%         \end{exerpart}  
% \end{exer}

% \begin{exer}[\exerE the Visitors do not fit in Hotel Infinity]
% \label{exer_hotel_visitors}
%     You are the proprietor of Hotel Infinity, as in \Cref{exer_hotel}. One day the Visitors arrive.  There is one visitor for every real number between 0 and 1.  Explain why you cannot fit all of the visitors in Hotel Infinity.  
% \end{exer}

% DO I WANT EXERCISES ON COUNTABLE??  Maybe a few

% %%%%%%%%%%%%%%%%%%%%%%%%%%%
% \subsection{\answers}
% %%%%%%%%%%%%%%%%%%%%%%%%%%%